\section{Divisibility at an inert prime}\label{sec:inert-divisibility}

Fix $p$ as in \S\ref{sec:inert-kernel}. Write $v_p(p)=1$,
$v_p(0)=+\infty$ and $\mathcal O=\mathcal O_{\mathbb C_p}$.
Choose $\varpi^{p^2-1}=p$ and put
\[
 \rho_p=\frac1{p^2-1},\qquad
 \gamma_p=\frac p{p^2-1}=\sum_{\substack{d\ge1\\d\ {\rm odd}}}p^{-d}.
\]
A derivative of order $D$ costs $\rho_pD$ on removing the scale
$\varpi$. The factorials in the derivative supply $D/(p-1)$ up to a
digit error. We compute how much of the difference survives at a top.

\subsection{The differential operators}

On one Frobenius orbit write $s\in\mathbb Z/2\nu\mathbb Z$ for its
positions. Let $X_s=(X_{s,a,b})$ be an $r_0\times r_1$ matrix and
$Y_s$ an $r_0\times r_0$ matrix of frame variables. For a sequence
$I$ of $h$ distinct column indices put
\[
 M_{s,I}=\det (Y_sX_s)_{[h],I},\qquad [h]=\{1,\ldots,h\}.
\]
Let $P$ be a product of such minors, one for each column of the weight
diagram at $s$, with $h$ its height. Fix a set $\mathcal T$ of tops.
Let $\alpha_s$ count the height-one columns at $s\in\mathcal T$, and
put $\alpha_s=0$ elsewhere. After removing these columns let $C_s$
be the number of remaining columns and $B_s$ their total height.
Thus $D=\deg_X P=\sum_s(\alpha_s+B_s)$. Write
$d_p(0)=0$ and $d_p(n)=1+\lfloor\log_p n\rfloor$ for $n\ge1$. Put
\[
 \Delta=\sum_s\{r_1d_p(\alpha_s)+r_0r_1d_p(B_s)\}
             +\sum_{s\in\mathcal T}d_p(\alpha_s)^2 .
\]

For $s\in\mathcal T$ let $E_s\subset\{1,3,\ldots,2\nu-1\}$ consist
of the least positive displacements from $s$ to tops of opposite parity.
Set $(z)_+=\max(z,0)$ and define
\begin{align}
 g_A&=\sum_{s\in\mathcal T}
       \sum_{\substack{d\ge1\ {\rm odd}\\d\notin E_s}}
       (\alpha_s p^{-d}-C_{s+d})_+-\Delta ,
       \label{eq:inert-base-gain}\\
 g_B&=\gamma_p\sum_{s\in\mathcal T_{\rm ev}}\alpha_s
       +\sum_{s\in\mathcal T_{\rm odd}}\ \sum_{d\ge1\ {\rm odd}}
       (\alpha_s p^{-d}-C_{s+d})_+-\Delta .
       \label{eq:inert-parity-gain}
\end{align}
The position $s+d$ is cyclic. The exclusion $d\in E_s$ uses the
integer $d$ itself.

For $J\subset[r_1]$ let $N_J$ count the removed columns with index
outside $J$ at even tops or inside $J$ at odd tops. Add to this
the sum of $|I\cap J|$ over the remaining minors at odd positions.
In particular $N_\varnothing=\sum_{s\in\mathcal T_{\rm ev}}\alpha_s$.

For coefficients $\upsilon_s\in\mathcal O$ put
\[
 S_{s,a,b}=\sum_{l\ge0}(-1)^l X_{s-2l,a,b}^{p^{2l}},\qquad
 h_\upsilon=\sum_{s,a,b}\upsilon_s S_{s,a,b}^{\,p}S_{s+1,a,b}.
\]
Write $\Theta_1\in\mathcal C(h)$ if
$\Theta_1=\sum_{i\ge0}p^{-i}h^i U_i$ with $U_i\in\mathcal O[[X]]$.
Here $h(0)=0$, so each degree uses finitely many $i$. We use the
orbit estimate of \citep[Theorem~9.6]{adriazola-grh} in the following form.

\begin{theorem}\label{thm:inert-valuation}
Suppose $\Theta(\varpi X)\in\mathcal C(h_\upsilon)$. Every coefficient
$c$ in the frame variables of $\langle P,\Theta\rangle
:=(P(\partial_X)\Theta)(0)$ satisfies
\begin{equation}\label{eq:inert-valuation}
       v_p(c)\ge g_A,\qquad v_p(c)+N_J\ge g_B
       \quad(J\subset[r_1]).
\end{equation}
\end{theorem}

Taking $J=\varnothing$ gives $v_p(c)\ge g_B-\sum_{s\ {\rm even}}\alpha_s$.
For a term carrying a factor $p^e$ with $e\ge N_J$, the second bound
gives $v_p(p^e c)\ge g_B$.

\begin{proof}
Keep $Y$ indeterminate and take all bounds coefficientwise in $Y$.
Write $\xi_{s,a}=(Y_s)_{1a}$ and
let $e_{s,b}$ count the removed columns at $s$ with index $b$.
Their factors in $P$ are $\sum_a\xi_{s,a}X_{s,a,b}$. Let $Q$ be the
product of the remaining minors. Set $t_{s,b}=0$ off the tops.
The chain rule gives
\[
 \varpi^D\langle P,\Theta\rangle
 =e!\,Q(\partial_X)[t^e]\,\Theta\bigl(\varpi(X+\xi t)\bigr)
       \big|_{X=0},
 \qquad (\xi t)_{s,a,b}=\xi_{s,a}t_{s,b}.
\]
For a multi-index $n$ we use $n!=\prod_i n_i!$.

\emph{Expansion modulo $p$.}
Modulo $p$, the terms of $X$-degree less than $p$ in
$h_\upsilon(X+\xi t)$ are integral linear combinations of
\begin{align*}
 &t_{s,b}^{p^d}\,t_{s+d-2k,b}^{p^{2k}},
                &&d\ge1\ {\rm odd},\quad k\ge0,\\
 &t_{s,b}^{p^d}\,L_{s,d,b}(X),\qquad
 L_{s,d,b}(X)=\sum_a\xi_{s,a}^{p^d}X_{s+d,a,b},
                &&d\ge1\ {\rm odd}.
\end{align*}
Indeed Frobenius distributes over sums modulo $p$. A first factor
$S^p$ contributing an $X$ has degree at least $p$. A second factor
$S$ contributing an $X$ has degree one only in its initial term, and
then the sum over $a$ is $L_{s,d,b}$.
Write $h'$ for the sum of these terms and an integral series $Z$ of
$X$-degree at least $p$, choosing lifts with zero constant term.
Then $h_\upsilon(X+\xi t)=h'+pV$ with $V$ integral and $V(0)=0$.
Expanding $(h'/p+V)^i$ absorbs $V$ into the integral factors.
We may truncate throughout at total degree $D$.

\emph{Factorials.}
In a term of $(h'/p)^i$ take $A$ pure $t$ factors, $B$ linear factors
and $z$ factors from $Z$, so $i=A+B+z$. Let $k_{s,d,b}$ be the
multiplicities of the linear factors and $b_{s,d}=\sum_b k_{s,d,b}$.
Move the linear factors onto $Q$ by differentiation. If
$\mathbf l_{s,d,b}$ is the coefficient vector of $L_{s,d,b}$, this
replaces $Q$ by
\[
 \Big(\prod_{s,d,b}k_{s,d,b}!\Big)
 [\tau^k]\,Q\Big(X+\sum_{s,d,b}\tau_{s,d,b}\mathbf l_{s,d,b}\Big).
\]
This follows from $\langle Q,LG\rangle=\langle D_{\mathbf l}Q,G\rangle$.
The multinomial coefficient times the displayed factorials is divisible
by $B!$, by first choosing the $B$ linear factors and then ordering
their types. For a surviving monomial $X^n$ of the translated
coefficient, $|n|=\sum_s B_s-B$, its degree at $s$ is at most $B_s$,
and $pz\le |n|$. Thus its contribution to the scaled pairing has
valuation at least
\[
 v_p(e!)+v_p(B!)+v_p(n!)-A-B-z .
\]
Legendre's formula gives $v_p(m!)\ge m/(p-1)-d_p(m)$.
At $s$ the entries of $e$ sum to $\alpha_s$, and $n$ has at most
$r_0r_1$ entries there. Also $\prod_{s,d} b_{s,d}!$ divides $B!$.
The available $t$-degree gives $p^d b_{s,d}\le\alpha_s$, so at each
top at most $d_p(\alpha_s)$ of these factorials are nontrivial and
each loses at most $d_p(\alpha_s)$. The three factorials therefore
give at least $D/(p-1)-\Delta$. Since $z\le\sum_s B_s/p$ and
$1/(p-1)-\rho_p=\gamma_p\ge1/p$, removing $\varpi^D$ leaves
\begin{equation}\label{eq:inert-term-bound}
              \gamma_p\sum_s\alpha_s-A-B-\Delta
\end{equation}
as a lower bound for each term.

\emph{Counting the factors.}
Assign each pure $t$ factor to the top carrying its larger exponent.
The exponents are $p^d$ and $p^{2k}$ with opposite parities of $d,2k$.
From the assigned top $s$ the other top lies a positive odd distance
$|d-2k|$ forward, modulo $2\nu$. Its least positive residue $j$
lies in $E_s$ and the exponent at $s$ is at least $p^j$.
Write $a_{s,j}$ for the resulting counts. For every positive odd $j$,
\[
   a_{s,j}+b_{s,j}\le \alpha_s p^{-j},\qquad
   a_{s,j}=0\ (j\notin E_s),\qquad b_{s,j}\le C_{s+j}.
\]
The first bound counts $t$-degree. For the last, hold $(s,j)$ fixed in
the translation of $Q$. At $s+j$ each changed column is a multiple of
the same vector $Y_{s+j}(\xi_{s,a}^{p^j})_a$. Multilinearity gives
degree at most one per minor in these translation variables, even with
the other translations present. Summing the three bounds, using
$\sum a_{s,j}=A$, $\sum b_{s,j}=B$ and $\sum_{j\ge1\ {\rm odd}}p^{-j}
=\gamma_p$, gives the first inequality in \eqref{eq:inert-valuation}
from \eqref{eq:inert-term-bound}.

For the second inequality retain the bounds
$b_{s,j}\le\min(\alpha_s p^{-j},C_{s+j})$ at odd tops. It remains to
pay for $A+\sum_{s\ {\rm even},j}b_{s,j}$. A pure factor uses one
index $b$ at both parities. It uses at least one counted $t$ at the
even top if $b\notin J$, and at the odd top if $b\in J$.
A linear factor at an even top with $b\notin J$ also uses a counted
$t$. Their total is bounded by the corresponding entries of $e$.
Each remaining linear factor, from an even top with $b\in J$, acts on
a column indexed by $b$ at an odd position. A minor there absorbs at
most $|I\cap J|$ such factors by multilinearity in those columns.
Hence $A+\sum_{s\ {\rm even},j}b_{s,j}\le N_J$.
Substitution in \eqref{eq:inert-term-bound} gives the second bound.
\end{proof}

\subsection{Application to the chosen shift}

Fix a lift $1\le j\le m$ and an auxiliary word $u$ at the shift $u_0$
of \S\ref{sec:theta-chain}. Take $r_0=N_{j-1}$ and $r_1=N_j$.
For a place $v$ of the maximal real subfield above $p$ split in $E$,
choose a place $w$ above $v$ and an embedding $E_w\hookrightarrow
W_p[1/p]$ inducing the fixed embedding on $K$. Write
$[E_w:K_p]=\nu$. The $2\nu$ positions index the embeddings
$\tau_s\in\Phi$ given by $\sigma^s$ on $w$ for even $s$ and by
$\sigma^s$ after conjugation for odd $s$.
Take $\mathcal T=\{s:\tau_s=(0,0,t),\ t\in T\}$.
At $s$ use the partition
\[
   \mu_s(i)=\lambda^{(j)}_{\tau_s}(i)+w_j(\tau_s).
\]
Its last $t_r+1$ entries vanish by \eqref{eq:chain-weights}, so
every column has height at most $r_0$. At a top,
$\alpha_s=\mu_s(1)-\mu_s(2)=a_t-\lambda^{(0)}_{\tau_s}(2)$
and $C_s=\mu_s(2)$. Elsewhere $C_s=\mu_s(1)$.
These are the shapes of the differential operators at both parities
\citep[Lemma~11.36]{adriazola-grh}.

The extension $E_w/K_p$ is unramified, so its integral trace pairing
is perfect. Choose dual bases $(\omega_l),(\omega_l^*)$ of
$\mathcal O_{E_w}$ over $\mathcal O_{K_p}$. In the hyperbolic frame
of \S\ref{sec:inert-kernel} put,
for each $a,b$,
\[
 L_p(x_{l,a,b})=\sum_{s\ {\rm even}}\sigma^s(\omega_l^*)X_{s,a,b},
 \qquad
 L_p(y_{l,a,b})=\sum_{s\ {\rm odd}}\sigma^s(\omega_l)X_{s,a,b}.
\]
These are the additive matrix coordinates in the two frames
\citep[Lemma~11.43]{adriazola-grh}.
Write $\vartheta_v=\exp(-\sum_{l,a,b}\beta(x_{l,a,b},y_{l,a,b}))$.
For every integral series $U_v$ in these formal parameters,
\begin{equation}\label{eq:inert-class}
         (\vartheta_v U_v)(\varpi X)\in\mathcal C(h_\upsilon)
\end{equation}
for some integral $\upsilon$. This is the height-two series calculation
of \citep[Proposition~9.5]{adriazola-grh}. Its hypotheses are precisely
the additivity and Frobenius relation of \S\ref{sec:inert-kernel}
and trace duality above. Indeed $[b]_F X=bX$ and $\sigma b=b^p$
for $b^{p^2-1}=1$. The series $\varpi^{-1}L_p(\varpi X)$ has integral
coefficients and linear coefficient one. Its compositional inverse is
integral, so the rescaled formal parameters preserve integrality of $U_v$.

\begin{corollary}\label{cor:inert-chain-jets}
Assume Hypothesis~\ref{hyp:inert-kernel-integrality}.
For each split $v$ choose a minor product $P_v$ as above and a set
$J_v\subset[r_1]$. In the additive coordinates, every frame coefficient
$c$ of $\langle\prod_v P_v,\Theta_{j,u,p}\rangle$ satisfies
\[
 v_p(c)\ge\sum_v\max\{g_{A,v},\,g_{B,v}-N_{J_v}\}.
\]
\end{corollary}

\begin{proof}
For the prescribed kernel in its prescribed frame, Hypothesis~\ref{hyp:inert-kernel-integrality}
gives the inclusion
\[
 U:=\Big(\prod_v\vartheta_v^{-1}\Big)\Theta_{j,u,p}
      \in\mathcal O[[\,x_{v,l,a,b},y_{v,l,a,b}\,]].
\]
The pairing uses only terms of $U$ of total degree at most
$\sum_v\deg P_v$, since the formal parameters have zero constant term.
Each such monomial is an integral coefficient times $\prod_v U_v$
with $U_v$ a monomial in the parameters of $v$. Its pairing factors as
$\prod_v\langle P_v,\vartheta_v U_v\rangle$. By
\eqref{eq:inert-class} and Theorem~\ref{thm:inert-valuation} the
valuations add across orbits. The ultrametric inequality gives the
same lower bound for the sum over monomials. With only Proposition~\ref{prop:inert-kernel},
the lower bound acquires the additional term $v_p(c_0)$ from its
undetermined scalar $c_0$.
\end{proof}
